% Standalone toggle: the parent paper defines \figisstandalone as 0
% BEFORE \input-ing this file, so \providecommand below is a no-op there.
% Compiled directly (e.g. in Overleaf on this file alone), \figisstandalone
% defaults to 1 and this file wraps itself in a full, compilable document.
\providecommand{\figisstandalone}{1}

\ifnum\figisstandalone=1
\documentclass[border=8mm]{standalone}
\usepackage{tikz}
\usetikzlibrary{positioning,arrows.meta,fit,backgrounds,calc,shapes.geometric}
\definecolor{pINK}{HTML}{3a352d}
\definecolor{pSUBINK}{HTML}{6f6757}
\definecolor{pCREAM}{HTML}{f4efe6}
\definecolor{pZERO}{HTML}{b8ad99}
\definecolor{pTERRA}{HTML}{c0603a}
\definecolor{pTEAL}{HTML}{3f6f74}
\definecolor{pSAGE}{HTML}{93a98d}
\definecolor{pSAND}{HTML}{cba14e}
\definecolor{pGRAY}{HTML}{a89f8d}

\begin{document}
\fi

% --- Custom TikZ Pic for 3D Isometric Cuboids ---
\tikzset{
    cuboid/.pic={
        \tikzset{/cuboid/.cd,#1}
        \pgfmathsetmacro{\cubew}{\pgfkeysvalueof{/cuboid/width}}
        \pgfmathsetmacro{\cubeh}{\pgfkeysvalueof{/cuboid/height}}
        \pgfmathsetmacro{\cubed}{\pgfkeysvalueof{/cuboid/depth}}
        \pgfmathsetmacro{\cubexy}{0.35}

        \coordinate (-left) at (0, 0);
        \coordinate (front-bl) at (0, -\cubeh/2);
        \coordinate (front-br) at (\cubew, -\cubeh/2);
        \coordinate (front-tr) at (\cubew, \cubeh/2);
        \coordinate (front-tl) at (0, \cubeh/2);
        \coordinate (back-bl) at (\cubed*\cubexy, -\cubeh/2 + \cubed*\cubexy);
        \coordinate (back-br) at (\cubew+\cubed*\cubexy, -\cubeh/2 + \cubed*\cubexy);
        \coordinate (back-tr) at (\cubew+\cubed*\cubexy, \cubeh/2 + \cubed*\cubexy);
        \coordinate (back-tl) at (\cubed*\cubexy, \cubeh/2 + \cubed*\cubexy);

        \draw[fill=\pgfkeysvalueof{/cuboid/fill}, draw=\pgfkeysvalueof{/cuboid/draw}, line width=0.7pt, join=round]
            (front-tl) -- (back-tl) -- (back-tr) -- (front-tr) -- cycle;
        \draw[fill=\pgfkeysvalueof{/cuboid/fill}, draw=\pgfkeysvalueof{/cuboid/draw}, line width=0.7pt, join=round]
            (front-br) -- (back-br) -- (back-tr) -- (front-tr) -- cycle;
        \draw[fill=\pgfkeysvalueof{/cuboid/fill}, draw=\pgfkeysvalueof{/cuboid/draw}, line width=0.7pt, join=round]
            (front-bl) -- (front-br) -- (front-tr) -- (front-tl) -- cycle;

        \coordinate (-right) at (\cubew, 0);
        \coordinate (-top) at (\cubew/2, \cubeh/2 + \cubed*\cubexy);
    }
}
\tikzset{
    /cuboid/.search also={/tikz},
    /cuboid/.cd,
    width/.initial=1.2, height/.initial=1.2, depth/.initial=1.2,
    fill/.initial=white, draw/.initial=black
}

\ifnum\figisstandalone=1\else
\begin{figure*}[t]
\centering
\resizebox{\textwidth}{!}{%
\fi
\begin{tikzpicture}[
  font=\sffamily\large,
  node distance=8mm,
  box/.style={draw, rounded corners=3pt, minimum height=12mm, minimum width=28mm,
              align=center, line width=0.7pt, text=pINK, inner sep=6pt},
  stage/.style={box, fill=pCREAM, draw=pZERO},
  fixed/.style={box, fill=pTEAL!18, draw=pTEAL},
  episodic/.style={box, fill=pTERRA!12, draw=pTERRA},
  oracle/.style={box, fill=pSAGE!25, draw=pSAGE},
  ourcode/.style={box, fill=pSAND!30, draw=pSAND},
  zsdata/.style={box, fill=pTERRA!12, draw=pTERRA, dash pattern=on 3pt off 2pt},
  panel/.style={draw=pGRAY, dashed, rounded corners=6pt, inner sep=14pt, line width=0.6pt},
  % tensor slivers: dashed border + lighter fill, to distinguish from solid module boxes
  sliver/.style={draw=pTEAL, fill=pTEAL!10, line width=0.7pt, minimum width=5mm,
                 inner sep=0pt, rounded corners=1.5pt, dash pattern=on 2pt off 1.5pt},
  sliver_ep/.style={draw=pTERRA, fill=pTERRA!10, line width=0.7pt, minimum width=5mm,
                    inner sep=0pt, rounded corners=1.5pt, dash pattern=on 2pt off 1.5pt},
  lbl/.style={text=pSUBINK, font=\sffamily\small},
  tinylbl/.style={text=pSUBINK, font=\sffamily\small},
  callout/.style={text=pSUBINK, font=\sffamily\small, align=left},
  panel_lbl/.style={text=pINK, font=\sffamily\large\bfseries},
  arr/.style={-{Triangle[length=2.5mm,width=2mm]}, line width=0.8pt, draw=pINK, rounded corners=3pt},
  loopback/.style={arr, draw=pTERRA},
  leader/.style={-{Triangle[length=2mm,width=1.5mm]}, line width=0.5pt, draw=pSUBINK, dash pattern=on 1.5pt off 1.5pt},
  zsarrow/.style={arr, draw=pTERRA, dash pattern=on 3pt off 2pt},
  legendkey/.style={draw, rounded corners=1.5pt, minimum width=5mm, minimum height=3mm, line width=0.6pt},
  legendkey_tensor/.style={draw, rounded corners=1.5pt, minimum width=5mm, minimum height=3mm, line width=0.6pt, dash pattern=on 2pt off 1.5pt},
]

% =========================================================
% GLOBAL GRID
% Two master rows drive the whole figure:
%   ROW_TOP  (y=0)    : netlist column, GCN encoder, pooled/curr slivers
%   ROW_BOT  (y=-46mm): occupancy grid, CNN encoder, cnn/wh slivers
% Columns are fixed x-coordinates so everything snaps to a grid.
% =========================================================
\def\rowtop{0}
\def\rowbot{-46mm}

% =========================================================
% PANEL A: INPUT ENCODING  (left column, two rows)
% =========================================================
% --- Top row: netlist pipeline ---
\node[episodic, minimum width=26mm] (netlist) at (-12mm,\rowtop) {Netlist\\[-0.3ex]\small(Verilog)};
\node[episodic, minimum width=26mm, right=11mm of netlist] (graph) {Netlist graph\\[-0.3ex]\small H-block nodes+fabric node};
\node[episodic, minimum width=26mm, right=11mm of graph] (pad) {Pad to\\[-0.3ex]fixed capacity\\[-0.3ex]\small max.\ nodes/edges};
\node[tinylbl, below=1mm of pad] (pad_tiny) {once per episode};

\node[callout, above=4mm of graph] (feat) {\texttt{[is\_dsp, is\_bram, is\_fabric, net\_count]}};
\draw[leader] (feat.south) -- (graph.north);

% --- Bottom row: occupancy grid, aligned under the graph column ---
\coordinate (occ_pos) at (graph |- {(0,-46mm)});
\node[minimum width=1cm, minimum height=1.6cm] (occ_anchor) at (occ_pos) {};
\pic (occ) at ($(occ_anchor.center) + (-0.6cm,-0.35cm)$) {cuboid={width=0.5, height=2.0, depth=2.0, fill=pTERRA!12, draw=pTERRA}};
\node[text=pINK, font=\sffamily\large\bfseries, align=center] (occ_lbl) at ($(occ_pos)+(0,-20mm)$)
      {Occupancy grid\\[-0.3ex]\normalfont\small current placement state, per step\\[-0.3ex]\normalfont\small once per step};
\node[minimum height=3mm, below=0mm of occ_lbl] (panelA_pad) {};

% --- Zero-shot unseen netlist, below netlist ---
\node[zsdata, minimum width=26mm] (unseen) at (netlist |- {(0,-32mm)}) {Unseen netlist\\[-0.3ex]\small never trained on};
\draw[zsarrow] (unseen.north) -- (netlist.south);
\node[lbl, below=2mm of unseen, text width=30mm, align=center] (unseen_lbl)
      {same checkpoint,\\[-0.3ex]single rollout $\to$ \S 5.3};

% --- valid(w,h): per-episode data, aligned under graph column ---
\node[episodic, text width=26mm, align=center] (validwh) at (graph |- {(0,-90mm)})
      {valid $(w,h)$\\[-0.3ex]\small norm.\ footprint};
% Invisible node to give Panel A a small right margin past validwh.
\node[minimum width=2mm, right=2mm of validwh] (validwh_pad) {};

\draw[arr] (netlist) -- (graph);
\draw[arr] (graph) -- (pad);

% =========================================================
% PANEL B: POLICY
% Encoders STACKED VERTICALLY: GCN on ROW_TOP, CNN on ROW_BOT, same column.
% =========================================================
\def\enccol{118mm}   % x-position of the encoder column

% --- GCN encoder (top row) ---
\node[fixed, minimum width=30mm, minimum height=26mm] (gcn_box) at (\enccol,\rowtop) {};
\node[text=pINK, font=\sffamily\large\bfseries] (gcn_lbl) at ($(gcn_box.north)+(0,5mm)$) {GCN encoder};
\node[text=pSUBINK, font=\sffamily\small, align=center] (gcn_sub) at ($(gcn_box.south)+(0,-5mm)$)
      {2$\times$ GCNConv (64)\\[-0.3ex]$\to$ node embeddings};
\begin{scope}[shift={($(gcn_box.center)+(-1.5mm, 0.5mm)$)}]
  \coordinate (n1) at (-7mm,-4mm); \coordinate (n2) at (-3mm,5mm);
  \coordinate (n3) at (6mm,6mm);   \coordinate (n4) at (7mm,-5mm);
  \coordinate (n5) at (1mm,-7mm);
  \draw[pTEAL, line width=1pt] (n1) -- (n2) -- (n3) -- (n4) -- (n5) -- (n1);
  \draw[pTEAL, line width=1pt] (n1) -- (n5) -- (n3); \draw[pTEAL, line width=1pt] (n2) -- (n5);
  \foreach \n in {n1,n2,n3,n4,n5} { \fill[white] (\n) circle (1.5mm); \draw[pTEAL, line width=1pt] (\n) circle (1.5mm); }
\end{scope}

% --- CNN encoder (bottom row, SAME column as GCN: vertically stacked) ---
\node[fixed, minimum width=30mm, minimum height=26mm] (cnn_box) at (\enccol,\rowbot) {};
\node[text=pINK, font=\sffamily\large\bfseries] (cnn_lbl) at ($(cnn_box.north)+(0,5mm)$) {CNN encoder};
\node[text=pSUBINK, font=\sffamily\small, align=center] (cnn_sub) at ($(cnn_box.south)+(0,-5mm)$)
      {4-channel grid\\[-0.3ex]AdaptiveAvgPool $\to$ Linear};
\begin{scope}[shift={($(cnn_box.center)+(-0.35,-0.05)$)}]
  \pic (cnn1) at (-1.0,-0.21)  {cuboid={width=0.6, height=1.3, depth=1.2, fill=pTEAL!18, draw=pTEAL}};
  \pic (cnn2) at (0.7,0)   {cuboid={width=0.8, height=0.75, depth=0.6, fill=pTEAL!18, draw=pTEAL}};
  
  \coordinate (s1) at (-0.16, 0.35);
  \coordinate (s2) at (-0.16, 0);
  \coordinate (s3) at (-0.16, -0.35);
  \coordinate (t1) at (0.7, 0.22);
  \coordinate (t2) at (0.7, 0);
  \coordinate (t3) at (0.7, -0.22);
  \coordinate (m1) at (0.35, 0.25);
  \coordinate (m2) at (0.35, 0);
  \coordinate (m3) at (0.35, -0.25);
  
  \foreach \src in {s1,s2,s3} {
    \foreach \mid in {m1,m2,m3} {
      \draw[pTEAL, opacity=0.35, line width=0.4pt] (\src) -- (\mid);
    }
  }
  \foreach \mid in {m1,m2,m3} {
    \foreach \tgt in {t1,t2,t3} {
      \draw[pTEAL, opacity=0.35, line width=0.4pt] (\mid) -- (\tgt);
    }
  }
  \foreach \mid in {m1,m2,m3} {
    \fill[white] (\mid) circle (1.0pt);
    \draw[pTEAL, line width=0.6pt] (\mid) circle (1.0pt);
  }
\end{scope}

% --- Encoder input arrows (clean horizontals) ---
\draw[arr] (pad.east) -- (gcn_box.west);
\draw[arr] ($(occ_pos) + (3mm,0)$) -- (cnn_box.west);

% =========================================================
% INTERMEDIATE EMBEDDING SLIVERS
% Stacked in one column to the right of the encoders.
% pool & curr come off GCN (top); cnn off CNN (bottom); wh off valid(w,h).
% =========================================================
\def\embcol{162mm}
% Two logical pairs: GCN outputs (top), then CNN + footprint (bottom).
% emb_cnn stays locked to the CNN box row (-46) so its arrow is dead horizontal.
\node[sliver, minimum height=12mm] (emb_pool) at (\embcol, 9mm) {};
\node[sliver, minimum height=12mm] (emb_curr) at (\embcol, -9mm) {};
\node[sliver, minimum height=12mm] (emb_cnn)  at (\embcol,-46mm) {};
\node[sliver_ep, minimum height=8mm] (emb_wh)  at (\embcol,-90mm) {};

% --- GCN fork: split node then route to pool & curr ---
\coordinate (gcn_split) at ($(gcn_box.east)+(8mm,0)$);

\draw[pINK, line width=0.8pt] (gcn_box.east) -- (gcn_split);
\draw[pINK, line width=0.8pt, rounded corners=3pt] (gcn_split) |- ($(emb_pool.west)+(-12mm,0)$) -- (emb_pool.west);
\draw[arr] ($(emb_pool.west)+(-2pt,0)$) -- (emb_pool.west);
\node[lbl, align=center, anchor=south] (pool_lbl)
      at ($(emb_pool.north) + (0, 0.5mm)$) {mean-pooled\\[-0.3ex]whole-netlist\\[-0.3ex]\small 64-d};

\draw[pINK, line width=0.8pt, rounded corners=3pt] (gcn_split) |- ($(emb_curr.west)+(-12mm,0)$) -- (emb_curr.west);
\draw[arr] ($(emb_curr.west)+(-2pt,0)$) -- (emb_curr.west);
\node[lbl, align=center, anchor=north] (curr_lbl)
      at ($(emb_curr.south) + (0, -0.5mm)$) {current\\[-0.3ex]H-block\\[-0.3ex]\small 64-d};

% --- CNN output: clean horizontal into emb_cnn (same row) ---
\draw[arr] (cnn_box.east) -- (emb_cnn.west);
\node[lbl, align=center, anchor=north] (cnn_lbl)
      at ($(emb_cnn.south) + (0, -0.5mm)$) {CNN output\\[-0.3ex]\small 64-d};

% --- valid(w,h) -> emb_wh : per-episode footprint crosses into the concat ---
\draw[arr] (validwh.east) -- (emb_wh.west);
\node[lbl, align=center, anchor=north] (wh_lbl)
      at ($(emb_wh.south) + (0, -0.5mm)$) {\small 2-d};

% =========================================================
% CONCATENATION SLIVER
% =========================================================
\def\concatcol{200mm}
\node[sliver, fill=pTEAL!22, minimum height=122mm] (concat) at (\concatcol,-40mm) {};
\node[text=pINK, font=\sffamily\large\bfseries, align=center] (concat_lbl) at ($(concat.north)+(0,7mm)$)
      {Concatenation\\[-0.3ex]\normalfont\small \textbf{194-d}};

\draw[arr] (emb_pool.east) -- (emb_pool.east -| concat.west);
\draw[arr] (emb_curr.east) -- (emb_curr.east -| concat.west);
\draw[arr] (emb_cnn.east)  -- (emb_cnn.east  -| concat.west);
\draw[arr] (emb_wh.east)   -- (emb_wh.east   -| concat.west);

% =========================================================
% FUSE  (a reduction: 194-d -> 128-d. Same sliver style as Concatenated,
% just shorter. No separate 128-d block — the Fuse output IS the 128-d.)
% =========================================================
\def\fusecol{234mm}
% Shorter sliver than concat (reduction); same thin-bar style.
\node[sliver, fill=pTEAL!22, minimum height=66mm] (fuse) at (\fusecol,-40mm) {};
\node[text=pINK, font=\sffamily\large\bfseries, align=center] (fuse_lbl) at ($(fuse.north)+(0,7mm)$)
      {Fuse \& Project\\[-0.3ex]\normalfont\small \textbf{128-d}};
\node[text=pSUBINK, font=\sffamily\small, align=center] (fuse_sub) at ($(fuse.south)+(0,-3.5mm)$)
      {Linear + ReLU};

% dense fan: concat right edge -> fuse left edge (the reduction)
\begin{scope}[on background layer]
  \foreach \i in {0.05,0.15,0.25,0.35,0.45,0.55,0.65,0.75,0.85,0.95}{
    \foreach \j in {0.15,0.4,0.6,0.85}{
      \draw[pTEAL, opacity=0.28, line width=0.5pt]
        ($(concat.south east)!\i!(concat.north east)$) --
        ($(fuse.south west)!\j!(fuse.north west)$);
    }
  }
\end{scope}
\foreach \i in {0.05,0.15,0.25,0.35,0.45,0.55,0.65,0.75,0.85,0.95}{
  \fill[pTEAL] ($(concat.south east)!\i!(concat.north east)$) circle (1.4pt);
}
\foreach \j in {0.15,0.4,0.6,0.85}{
  \fill[pTEAL] ($(fuse.south west)!\j!(fuse.north west)$) circle (1.4pt);
}
% Direction indicator: the fan flows concat -> fuse (the reduction).
\draw[arr, line width=1pt] ($(concat.east)!0.5!(fuse.west) + (-4mm,0)$) -- ($(concat.east)!0.5!(fuse.west) + (4mm,0)$);

% =========================================================
% PPO AGENT
% =========================================================
\node[fixed, minimum height=11mm, minimum width=24mm] (actor)  at ($(fuse.east)+(30mm,8mm)$)  {Actor\\[-0.3ex]\small (masked action)};
\node[fixed, minimum height=11mm, minimum width=24mm] (critic) at ($(fuse.east)+(30mm,-12mm)$) {Critic\\[-0.3ex]\small (value)};
\node[draw=pTEAL, dashed, rounded corners=5pt, fit=(actor)(critic), inner sep=7pt] (ppo_sys) {};
\node[text=pINK, font=\sffamily\large\bfseries] (ppo_lbl) at ($(ppo_sys.north)+(0,5mm)$) {PPO Agent};
\node[lbl, align=center, text width=38mm] (mask_lbl) at ($(ppo_sys.south)+(0,-5mm)$)
      {mask: valid coords,\\[-0.3ex]active benchmark\\[-0.3ex]loop: repeats per H-block};

% fuse -> ppo agent container
\draw[arr] (fuse.east) -- (fuse.east -| ppo_sys.west);



% =========================================================
% PANEL C: EVALUATION ORACLE
% =========================================================
\node[ourcode, minimum width=30mm, right=30mm of actor] (bake) {Render layout\\[-0.3ex]\small Jinja2 $\to$ arch.\ XML};
\node[oracle, minimum width=30mm, right=12mm of bake] (vtr)
      {VTR oracle\\[-0.3ex]\small synth/pack/place/route\\[-0.3ex]\small one call per episode};
\node[oracle, minimum width=30mm, below=12mm of vtr] (reward) {ADP $\to$ reward\\[-0.3ex]\small Eq. 3};
\node[ourcode, minimum width=30mm, below=12mm of bake] (baseline)
      {Baseline metrics\\[-0.3ex]\small Area$_0$,Delay$_0$,Power$_0$};

\draw[arr] (actor.east) -- (bake.west);
\draw[arr] (bake) -- (vtr);
\draw[arr] (vtr) -- (reward);
\draw[arr] (baseline) -- (reward);

% =========================================================
% PANELS (fit to content) + LOOPBACK
% =========================================================
\begin{scope}[on background layer]
\node[panel,
      fit=(netlist)(graph)(pad)(pad_tiny)(feat)(occ_anchor)(occ_lbl)(panelA_pad)(unseen)(unseen_lbl)(validwh)(validwh_pad),
      label={[panel_lbl]above:Input Encoding (\S 3.3)}] (panelA) {};
\node[panel,
      fit=(gcn_box)(gcn_lbl)(gcn_sub)(cnn_box)(cnn_lbl)(cnn_sub)
          (emb_pool)(emb_curr)(emb_cnn)(emb_wh)(concat)(concat_lbl)
          (fuse)(fuse_lbl)(fuse_sub)(ppo_sys)(ppo_lbl)(mask_lbl),
      label={[panel_lbl]above:Policy, fixed weights, trained end-to-end (\S 3.2, \S 3.3)}] (panelB) {};
\node[panel,
      fit=(bake)(vtr)(reward)(baseline),
      label={[panel_lbl]above:Evaluation Oracle (\S 4.1)}] (panelC) {};
\end{scope}

\coordinate (loopback_y) at (0,-120mm);
\draw[loopback] (panelC.south) -- (panelC.south |- loopback_y)
    -| node[pos=0.20, above=1.5mm, lbl, text width=46mm, align=center]
       {PPO gradient (training only)\\[-0.3ex]\small updates GCN, CNN, Fuse, Actor, Critic}
       (panelB.south);

% =========================================================
% LEGEND
% =========================================================
\node[legendkey, fill=pTERRA!12, draw=pTERRA] (lg1) at ($(current bounding box.south)+(-162mm,-16mm)$) {};
\node[lbl, anchor=west, right=2.5mm of lg1] (lg1t) {per-episode data};
\node[legendkey, fill=pTEAL!18, draw=pTEAL, right=8mm of lg1t] (lg2) {};
\node[lbl, anchor=west, right=2.5mm of lg2] (lg2t) {fixed policy module};
\node[legendkey_tensor, fill=pTEAL!10, draw=pTEAL, right=8mm of lg2t] (lg2c) {};
\node[lbl, anchor=west, right=2.5mm of lg2c] (lg2ct) {activation tensor};
\node[legendkey, fill=pSAND!30, draw=pSAND, right=8mm of lg2ct] (lg2b) {};
\node[lbl, anchor=west, right=2.5mm of lg2b] (lg2bt) {our code (rendering)};
\node[legendkey, fill=pSAGE!25, draw=pSAGE, right=8mm of lg2bt] (lg3) {};
\node[lbl, anchor=west, right=2.5mm of lg3] (lg3t) {external oracle (VTR)};
\node[legendkey, fill=pTERRA!12, draw=pTERRA, dash pattern=on 3pt off 2pt, right=8mm of lg3t] (lg4) {};
\node[lbl, anchor=west, right=2.5mm of lg4] (lg4t) {zero-shot data/path};
\node[minimum width=5mm, minimum height=3mm, right=8mm of lg4t] (lg5) {\tikz{\draw[loopback] (0,0) -- (5mm,0);}};
\node[lbl, anchor=west, right=2.5mm of lg5] (lg5t) {training-only update};

\end{tikzpicture}%
\ifnum\figisstandalone=1\else
}
\caption{System overview of the RL-based fabric layout search. Panel A encodes the input netlist and grid. Panel B rollouts the placement sequentially, positioning one H-block per step under action masking. Panel C uses Jinja2 to render the final placement into a VTR architecture XML for verification and reward computation. The dashed line (bottom-left) denotes the zero-shot path where unseen netlists bypass training and are directly evaluated by the fixed policy.}
\label{fig:architecture}
\end{figure*}
\fi

\ifnum\figisstandalone=1
\end{document}
\fi